% ============================================================================
% Primer avance del proyecto — Curso de Big Data
% Formato IEEE (conference)
% ============================================================================
\documentclass[conference]{IEEEtran}

\usepackage[utf8]{inputenc}
\usepackage[T1]{fontenc}
\usepackage[spanish]{babel}
\usepackage{amsmath,amssymb}
\usepackage{booktabs}
\usepackage{array}
\usepackage{graphicx}
\usepackage[hidelinks]{hyperref}

% ----------------------------------------------------------------------------
% Metadatos
% ----------------------------------------------------------------------------
\newcommand{\tituloProyecto}{Análisis de datos de telemetría de Fórmula 1:
caracterización y propuesta preliminar de resultados}
\newcommand{\autores}{[Nombre del estudiante]}
\newcommand{\institucion}{Universidad [Nombre]\\Facultad de [Ingeniería]}

\begin{document}

\title{\tituloProyecto}

\author{\IEEEauthorblockN{\autores}
\IEEEauthorblockA{\institucion}}

\markboth{Primer avance del proyecto — Curso de Big Data}{}%

\maketitle

% ----------------------------------------------------------------------------
% Resumen
% ----------------------------------------------------------------------------
\begin{abstract}
Se presenta el primer avance del proyecto del curso de Big Data, cuyo objetivo
es el análisis de un repositorio de datos de telemetría pública de Fórmula 1.
El repositorio seleccionado reúne los conjuntos de datos de TracingInsights
para las temporadas 2021 a 2025, organizados en archivos JSON de alta
frecuencia por temporada, Gran Premio, sesión, piloto y vuelta. Se describen
en detalle la estructura, los tipos de archivo y las variables disponibles,
así como el volumen aproximado de la información. Finalmente, se plantea una
propuesta preliminar de resultados y análisis, susceptible de modificarse
conforme avance el proyecto.
\end{abstract}

\begin{IEEEkeywords}
Big Data, Fórmula 1, telemetría, JSON, análisis de datos, TracingInsights.
\end{IEEEkeywords}

% ----------------------------------------------------------------------------
% 1. Definición del repositorio seleccionado
% ----------------------------------------------------------------------------
\section{Definición del repositorio seleccionado}

El repositorio seleccionado es una colección local de los conjuntos de datos
públicos de telemetría de Fórmula 1 publicados por \textbf{TracingInsights}
(\url{https://tracinginsights.com/data/}), correspondientes a las temporadas
\textbf{2021--2025}. Se trata de un repositorio no oficial, no afiliado a la
Fórmula 1, a la FIA ni a ningún equipo, distribuido bajo licencia
\textbf{Apache-2.0}.

El repositorio está organizado de forma jerárquica y autocontenida por
temporada. Cada carpeta anual (\texttt{2021/}, \texttt{2022/}, \texttt{2023/},
\texttt{2024/} y \texttt{2025/}) replica la estructura del repositorio GitHub
de TracingInsights correspondiente a ese año. La información se almacena
principalmente en archivos \textbf{JSON} y se organiza de acuerdo con la
siguiente jerarquía:

\begin{itemize}
  \item \textbf{Temporada} (por ejemplo, \texttt{2024/}).
  \item \textbf{Evento o Gran Premio} (por ejemplo,
        \texttt{Bahrain Grand Prix/}).
  \item \textbf{Sesión} (por ejemplo, \texttt{Race/}, \texttt{Qualifying/},
        \texttt{Practice 1/}).
  \item \textbf{Piloto} (por ejemplo, \texttt{VER/}, \texttt{HAM/}).
  \item \textbf{Vuelta} (por ejemplo, \texttt{1\_tel.json}).
\end{itemize}

Además de los datos, la colección incluye scripts de extracción y análisis en
Python, caches de FastF1 y metadatos de Git. Cabe señalar que los directorios
\texttt{.git/} contienen el historial completo de cada repositorio de temporada
y dominan la huella en disco, aunque no forman parte del conjunto de datos
utilizable.

En la Tabla~\ref{tab:resumen} se resume el volumen de información por
temporada.

\begin{table}[!htbp]
\centering
\caption{Resumen de almacenamiento por temporada.}
\label{tab:resumen}
\begin{tabular}{@{}cccc@{}}
\toprule
\textbf{Año} & \textbf{GPs} & \textbf{Archivos JSON} & \textbf{Tamaño JSON} \\
\midrule
2021 & 22 & 62{,}996 & 11.27 GB \\
2022 & 22 & 62{,}999 & 10.21 GB \\
2023 & 22 & 62{,}608 & 10.18 GB \\
2024 & 23 & 70{,}668 & 11.65 GB \\
2025 & 23 & 106{,}663 & 13.74 GB \\
\bottomrule
\end{tabular}
\end{table}

% ----------------------------------------------------------------------------
% 2. Descripción detallada de los datos
% ----------------------------------------------------------------------------
\section{Descripción detallada de los datos que se van a analizar}

El conjunto de datos se compone de varios tipos de archivos JSON, cada uno con
un propósito específico dentro de una sesión de Gran Premio. A continuación se
describen los más relevantes para el análisis.

\subsection{Telemetría por vuelta (\texttt{\{n\}\_tel.json})}

Se almacena dentro de cada carpeta de piloto y contiene muestras de telemetría
de alta frecuencia para una única vuelta. Cada archivo agrupa arreglos
paralelos (una entrada por muestra) con las siguientes variables principales:

\begin{itemize}
  \item \textbf{time}: tiempo transcurrido en la vuelta (segundos).
  \item \textbf{speed}: velocidad del monoplaza (km/h).
  \item \textbf{throttle}: porcentaje de presión del acelerador.
  \item \textbf{brake}: aplicación del freno.
  \item \textbf{gear}: marcha seleccionada.
  \item \textbf{rpm}: revoluciones del motor.
  \item \textbf{drs}: estado o zona de activación del DRS.
  \item \textbf{acc\_x}, \textbf{acc\_y}, \textbf{acc\_z}: aceleraciones
        lateral, longitudinal y vertical (en g).
  \item \textbf{x}, \textbf{y}, \textbf{z}: coordenadas de posición en pista.
  \item \textbf{distance} y \textbf{rel\_distance}: distancia dentro de la
        vuelta (metros) y distancia relativa.
  \item \textbf{DistanceToDriverAhead}: distancia al monoplaza que precede
        (metros).
  \item \textbf{DriverAhead}: código del piloto que precede.
\end{itemize}

Esta es la fuente de datos más densa del repositorio: una sola sesión de
carrera de un piloto puede contener más de 50 archivos de telemetría por
vuelta, y cada archivo miles de muestras de alta frecuencia.

\subsection{Tiempos por vuelta (\texttt{laptimes.json})}

Contiene los tiempos resumidos de toda la sesión para un piloto. Además de la
lista de tiempos y números de vuelta, incluye:

\begin{itemize}
  \item \textbf{compound}: compuesto del neumático utilizado (SOFT, MEDIUM,
        HARD, etc.).
  \item \textbf{stint}: número de stint (periodo entre paradas en boxes).
  \item \textbf{s1}, \textbf{s2}, \textbf{s3}: tiempos de los sectores 1, 2 y
        3.
  \item Variables adicionales como posición, estado de la vuelta, vida del
        neumático (\texttt{life}) y bandera de vuelta rápida personal
        (\texttt{pb}).
\end{itemize}

\subsection{Tiempos por vuelta de sesión (\texttt{session\_laptimes.json})}

Se almacena en la raíz de la sesión y consolida en un único archivo los arreglos
de \texttt{laptimes.json} de todos los pilotos, facilitando los análisis
comparativos a nivel de sesión.

\subsection{Metadatos de pilotos (\texttt{drivers.json})}

Contiene la lista de pilotos participantes con su código, nombre, apellido,
equipo, número de monoplaza, color del equipo y URL de la imagen oficial.

\subsection{Datos meteorológicos (\texttt{weather.json})}

Contiene muestras de clima en pista, entre ellas la temperatura ambiente
(\texttt{wAT}), la humedad (\texttt{wH}), la presión (\texttt{wP}), la
temperatura de pista (\texttt{wTT}), la dirección (\texttt{wWD}) y velocidad
(\texttt{wWS}) del viento, y un indicador de lluvia (\texttt{wR}).

\subsection{Datos de curvas (\texttt{corners.json})}

Describe la información geométrica de las curvas del circuito: número de
curva, coordenadas, ángulo, distancia dentro de la vuelta y rotación.

\subsection{Mensajes de dirección de carrera (\texttt{rcm.json})}

Contiene los mensajes emitidos por la dirección de carrera durante la sesión
(por ejemplo, sanciones o eliminaciones de vuelta por límites de pista).

\subsection{Clasificaciones (\texttt{standings.json}) — solo 2025}

Contiene la clasificación del campeonato de pilotos con posición, puntos,
victorias, podios, equipo, nacionalidad y diversas métricas de rendimiento.
Según la documentación del repositorio, este archivo puede estar desactualizado
(reporta la temporada 2023), por lo que deberá validarse antes de su uso.

En términos globales, el conjunto de datos utilizable está compuesto por
aproximadamente \textbf{367{,}000 archivos JSON} con un tamaño cercano a
\textbf{57 GB}, lo que lo convierte en un caso de estudio representativo de un
volumen de datos de tipo \emph{big data} dentro del dominio del deporte motor.

% ----------------------------------------------------------------------------
% 3. Propuesta preliminar de resultados
% ----------------------------------------------------------------------------
\section{Propuesta preliminar de resultados}

La siguiente propuesta de resultados es preliminar y está sujeta a cambios
conforme avance el proyecto. En esta primera etapa se plantean los siguientes
análisis y entregables:

\subsection{Análisis comparativo del rendimiento entre pilotos y equipos}
Comparar el rendimiento de los pilotos y equipos a partir de las telemetrías y
los tiempos por vuelta, identificando diferencias en sectores específicos
(S1, S2, S3) y en variables como velocidad, frenado, aceleración y uso del DRS.

\subsection{Análisis de estrategia de neumáticos y stints}
Estudiar el impacto del compuesto del neumático y de la longitud de cada stint
en los tiempos de vuelta, cuantificando la degradación del neumático a lo largo
de su vida útil.

\subsection{Impacto del clima en el rendimiento}
Correlacionar las variables meteorológicas (temperatura de pista, humedad,
velocidad del viento, lluvia) con los tiempos de vuelta y la telemetría, con el
fin de cuantificar su influencia en el rendimiento.

\subsection{Detección y análisis de adelantamientos}
Identificar maniobras de adelantamiento mediante la variable
\texttt{DistanceToDriverAhead} y el código del piloto que precede, y
caracterizar las zonas del circuito donde ocurren con mayor frecuencia.

\subsection{Evolución del rendimiento entre temporadas (2021--2025)}
Analizar la evolución del rendimiento a lo largo de las cinco temporadas
disponibles, comparando tiempos de vuelta, velocidades punta y consistencia
entre circuitos que se repiten año tras año.

\subsection{Resultados esperados}
Como resultado preliminar se espera generar reportes estadísticos y
visualizaciones (gráficas de telemetría, comparativas por sector, mapas de
degradación de neumáticos y distribuciones de tiempos) que permitan
caracterizar el rendimiento en Fórmula 1. Los hallazgos y el alcance de estos
análisis podrán ajustarse en función de la calidad de los datos y de los
recursos computacionales disponibles.

% ----------------------------------------------------------------------------
% Referencias
% ----------------------------------------------------------------------------
\begin{thebibliography}{1}
\bibitem{tracinginsights}
TracingInsights, ``F1 telemetry datasets,'' 2025. [En línea]. Disponible:
\url{https://tracinginsights.com/data/}.
\bibitem{fastf1}
theOehrly, ``FastF1: A python package for accessing and analyzing Formula 1
results, schedules, timing data and telemetry.'' [En línea]. Disponible:
\url{https://github.com/theOehrly/FastF1}.
\bibitem{jolpica}
Jolpica, ``Jolpica F1 (Ergast) — An open F1 API.'' [En línea]. Disponible:
\url{https://github.com/jolpica/jolpica-f1}.
\end{thebibliography}

\end{document}
